%! Author = Admin
%! Date = 18.04.2026

\section{Тестирование корректной работы контроллера}

Разработана автоматизированная проверка корректности выполнения симуляции.
Каждый тест запускает полное моделирование движения квадрокоптера по одной из типовых кривых и проверяет финальную норму боковой ошибки
\begin{gather}
    \left\|
        \begin{bmatrix}
            e_1(T) \\
            e_2(T)
        \end{bmatrix}
    \right\| < \varepsilon_{\mathrm{pass}},
    \qquad
    \varepsilon_{\mathrm{pass}} = 1.5 \ \text{м}.
\end{gather}
Набор тестовых траекторий включает:
\begin{enumerate}
    \item круговую спираль с радиусом $r=3$;
    \item горизонтальную окружность радиуса $r=3$;
    \item спираль с радиусом $r=2$;
    \item пространственную прямую $x=s,\ y=s,\ z=s$;
    \item прямую в плоскости $XZ$: $x=s,\ y=0,\ z=0.5s$;
    \item плотную спираль с радиусом $r=1.5$.
\end{enumerate}\hfill\break
Для прямых используется аналитическая функция ближайшей точки, а для спиралей - динамический наблюдатель.
Такой подход позволяет проверить корректность регулятора при разных значениях кривизны, различной норме касательного вектора и разных начальных условиях.\\
В ходе тестирования анализируются:
\begin{enumerate}
    \item финальные ошибки $e_1(T)$ и $e_2(T)$;
    \item норма скорости $\|v(T)\|$;
    \item устойчивость наблюдателя ближайшей точки;
    \item корректность выбора шага $\Delta t$ при заданном $\kappa$.
\end{enumerate}\hfill\break
Кроме логической проверки \texttt{PASS/FAIL}, для каждого сценария автоматически сохраняются графики траектории, ошибок, скорости и углов.
Это упрощает локализацию режимов, в которых требуется перенастройка коэффициентов.

На рисунках~\ref{fig:line_3d}--\ref{fig:circle_errors} приведены результаты трёх типовых сценариев: движение по пространственной прямой, спирали и горизонтальной окружности.

\begin{figure}[htbp]
    \centering
    \begin{subfigure}[b]{0.48\textwidth}
        \centering
        \includegraphics[width=\textwidth]{images/ch4_line_3d}
        \caption{Траектория, прямая}
        \label{fig:line_3d}
    \end{subfigure}
    \hfill
    \begin{subfigure}[b]{0.48\textwidth}
        \centering
        \includegraphics[width=\textwidth]{images/ch4_line_errors}
        \caption{Ошибки слежения, прямая}
        \label{fig:line_errors}
    \end{subfigure}
    \caption{Результаты моделирования: пространственная прямая ($V^{\ast}=1.0$, $T=40$ с)}
    \label{fig:line_results}
\end{figure}

\begin{figure}[htbp]
    \centering
    \begin{subfigure}[b]{0.48\textwidth}
        \centering
        \includegraphics[width=\textwidth]{images/ch4_spiral_3d}
        \caption{Траектория, спираль $r{=}3$}
        \label{fig:spiral_3d}
    \end{subfigure}
    \hfill
    \begin{subfigure}[b]{0.48\textwidth}
        \centering
        \includegraphics[width=\textwidth]{images/ch4_spiral_errors}
        \caption{Ошибки слежения, спираль $r{=}3$}
        \label{fig:spiral_errors}
    \end{subfigure}
    \caption{Результаты моделирования: спираль $r{=}3$ ($V^{\ast}=1.0$, $T=40$ с)}
    \label{fig:spiral_results}
\end{figure}

\begin{figure}[htbp]
    \centering
    \begin{subfigure}[b]{0.48\textwidth}
        \centering
        \includegraphics[width=\textwidth]{images/circle_r3_z5_traj_3d}
        \caption{Траектория, окружность $r{=}3$, $z{=}5$}
        \label{fig:circle_3d}
    \end{subfigure}
    \hfill
    \begin{subfigure}[b]{0.48\textwidth}
        \centering
        \includegraphics[width=\textwidth]{images/circle_r3_z5_errors}
        \caption{Ошибки слежения, окружность}
        \label{fig:circle_errors}
    \end{subfigure}
    \caption{Результаты моделирования: горизонтальная окружность $r{=}3$, $z{=}5$ ($V^{\ast}=1.0$, $T=40$ с)}
    \label{fig:circle_results}
\end{figure}